% Generated draft: requires mathematical and source review.
\documentclass[11pt]{article}
\usepackage[margin=25mm]{geometry}
\usepackage{fontspec}
\setmainfont{Latin Modern Roman}
\usepackage{amsmath,amssymb,mathtools,unicode-math}
\setmathfont{Latin Modern Math}
\usepackage{enumitem,needspace,xurl,hyperref}
\setlength{\parindent}{0pt}
\setlength{\parskip}{5pt}
\setlength{\emergencystretch}{3em}
\newcommand{\N}{\mathbb{N}}
\newcommand{\Z}{\mathbb{Z}}
\newcommand{\Q}{\mathbb{Q}}
\newcommand{\R}{\mathbb{R}}
\newcommand{\C}{\mathbb{C}}
\newcommand{\F}{\mathbb{F}}
\newcommand{\A}{\mathbb{A}}
\newcommand{\PP}{\mathbb{P}}
\newcommand{\E}{\mathbb{E}}
\newcommand{\eps}{\varepsilon}
\newcommand{\dd}{\,\mathrm{d}}
\begin{document}
\section*{UnsolvedMath ID: 30002579}
\textbf{Source problem code:} catalogue entry
\subsection{Definitions and mathematical statement}
The supplied record bundles several questions about Hopf actions and homological properties. Their standing hypotheses are incomplete, so the formulations below retain the unresolved parameters explicitly.

Let \(k\) denote the unspecified ground field, and let \(H\) be a Hopf algebra over \(k\), with coproduct \(\Delta\), counit \(\varepsilon\), and antipode \(S\). A left \(H\)-module algebra is a unital associative \(k\)-algebra \(A\) with an \(H\)-action satisfying
\[
 h\cdot(ab)=\sum (h_{(1)}\cdot a)(h_{(2)}\cdot b),\qquad h\cdot1_A=\varepsilon(h)1_A,
\]
where \(\Delta(h)=\sum h_{(1)}\otimes h_{(2)}\). A two-sided ideal \(I\subseteq A\) is \(H\)-stable if \(H\cdot I\subseteq I\). The algebra \(A\) is \(H\)-semiprime if it has no nonzero nilpotent \(H\)-stable two-sided ideal, and \(H\)-simple if \(A\ne0\) and its only \(H\)-stable two-sided ideals are \(0\) and \(A\). Ordinary semiprimeness excludes all nonzero nilpotent two-sided ideals. Write
\[
 Z(A)=\{z\in A:za=az\text{ for every }a\in A\}.
\]
Being finite over the center means being finitely generated as a \(Z(A)\)-module. Right noetherian means that ascending chains of right ideals stabilize; right artinian means that descending chains of right ideals stabilize.

\begin{enumerate}
\item For every \(H\)-semiprime right-noetherian \(H\)-module algebra \(A\) in the intended class, does the classical right quotient ring \(Q_{\mathrm{cl}}^r(A)\) exist and is it artinian? Here the classical right quotient ring is the right Ore localization at the set \(C\) of regular elements, namely elements that are neither left nor right zero divisors. In particular, its existence requires the corresponding right denominator conditions, including that for every \(a\in A\) and \(s\in C\) there are \(a'\in A\) and \(s'\in C\) with \(as'=sa'\). The record does not specify whether ``artinian'' requires both sides or only the right side, nor does it supply further restrictions on \(H\).

\item Does there exist a suitable \(H\)-simple, semiprime, noetherian \(H\)-module algebra \(A\), finite over \(Z(A)\), with infinite global dimension? Left global dimension is the supremum of the projective dimensions of all left modules, and right global dimension is defined analogously. The required side of global dimension and of noetherianity, and the meaning of ``suitable'' in the original statement, are unspecified.

\item Does there exist a suitable \(H\)-simple noetherian \(H\)-module algebra \(A\), finite over \(Z(A)\), that is not Gorenstein? The original statement does not include ordinary semiprimeness in this question, whereas the normalized statement groups it with the preceding question under a semiprimeness hypothesis. These two possible scopes must be distinguished. The meaning of Gorenstein is also unspecified: finite left and right self-injective dimensions, graded Artin--Schelter Gorenstein conditions, and other noncommutative Gorenstein conventions cannot be silently identified.

\item Suppose that \(H\) is cosemisimple and acts inner faithfully on a polynomial-identity domain \(A\) of degree \(d\). Must
\[
 \operatorname{PIdeg}(H^*)\le d^2
\]
hold? Cosemisimplicity means that the coalgebra underlying \(H\) is a direct sum of simple coalgebras. Inner faithfulness means that no nonzero Hopf ideal \(J\subseteq H\) annihilates \(A\), that is, \(J\cdot A=0\) implies \(J=0\). A polynomial-identity algebra satisfies a nonzero polynomial identity in noncommuting variables. Under the usual central-quotient convention for a PI domain, degree \(d\) means
\[
 \dim_{Z(Q)}Q=d^2,
\]
where \(Q\) is its classical division ring of fractions. The supplied record does not specify the meaning of \(H^*\), hypotheses making it an algebra with a defined PI degree, or the convention for that degree. The displayed inequality therefore remains conditional on these missing specifications.

\item For every noetherian Artin--Schelter regular algebra \(A\), does its given multiplication admit a braided Hopf algebra structure? A braided Hopf algebra is a Hopf algebra object in a specified braided monoidal category: in addition to multiplication and unit, it has a coassociative coproduct, counit, and antipode, with compatibility using the category's braiding. The ambient category, allowed braiding, grading requirements, and compatibility with the given algebra are not supplied. One common connected graded convention for Artin--Schelter regularity requires \(A=\bigoplus_{n\ge0}A_n\), \(A_0=k\), local finiteness, finite global dimension \(n\), finite Gelfand--Kirillov dimension, and left and right Artin--Schelter Gorenstein conditions. The latter require the corresponding \(\operatorname{Ext}^i(k,A)\) groups to vanish for \(i\ne n\) and to be a grading shift of the trivial module for \(i=n\). This convention is explanatory only: the record does not identify which definition of Artin--Schelter regularity is intended.
\end{enumerate}

No characteristic, algebraic-closure assumption, finite-dimensionality assumption on \(H\), or bijectivity assumption on its antipode is supplied. The intended quantification over fields and Hopf algebras consequently remains unresolved. The dated background assessment also mentions orbit questions but provides no orbit problem to formulate.
\subsection{Short English statement}
For permanent UnsolvedMath ID \(30002579\), the record asks about artinian quotient rings of Hopf-module algebras, examples with infinite global dimension or failure of Gorensteinness, a proposed PI-degree bound for a dual Hopf algebra, and braided Hopf structures on Artin--Schelter regular algebras. Missing standing hypotheses, conflicting semiprimeness assumptions, and unspecified dual, Gorenstein, and braided-category conventions prevent a fully determined statement. The supplied open status and verification label require independent human review.
\subsection{Sources}
\begin{itemize}
\item Local registry: \texttt{data/unsolvedmath/problems.json}, permanent UnsolvedMath ID \(30002579\), source problem code \texttt{catalogue entry}, title ``Quotient Rings and Homological Properties of Hopf-Module Algebras.'' The supplied record belongs to the Oberwolfach Reports Open Problems collection. No individual proposer is supplied; the proposed year is recorded as 2014.
\item Supplied source citation: \emph{Mini-Workshop: Infinite Dimensional Hopf Algebras (2015)}, Oberwolfach Reports, pp.\ 1111--1137, DOI \texttt{10.4171/owr/2014/20}. Source URL and literature-source label ``Oberwolfach open-questions report'': \url{https://doi.org/10.4171/owr/2014/20}. The citation's 2015 label and the record's 2014 year information are retained without reconciliation.
\item Supplied literature source, ``Hopf algebra orbits'': \url{https://dspace.kpfu.ru/xmlui/handle/net/139167}. No authors, theorem numbers, or detailed bibliographic citation are supplied.
\item Supplied literature source, ``Hopf-module-algebra Gorensteinness'': \url{https://doi.org/10.1017/S0017089516000185}. No authors or theorem numbers are supplied.
\item Literature assessment checked at \texttt{2026-08-21}: this is supplied dated curation, not independently verified. It labels the record open, states that theorems settle several restricted Hopf-action situations, including module-finite-over-center cases, and describes broader quotient-ring, orbit, homological, Gorenstein, PI-degree, and braided-Hopf questions as unresolved in general. The background explicitly requests independent verification of time-sensitive status and cited claims.
\item The supplied statement-status label is \texttt{corrected\_verified}; its verification note says that mathematical content was unchanged and notation, typography, and wording were normalized. This is a registry assertion, not independent mathematical verification of this draft.
\end{itemize}
\end{document}
